%!TEX root=../../main.tex

\limitingActivationSet*
\begin{proof}
    We say \( \Delta \preceq \Delta' \) if every entry of
    \( \Delta' \) is at least as large as \( \Delta \).
    If \( \Delta \prec \Delta' \), then it is straightforward to embed
    any network weights \( \theta \) conformal to \( \Delta \) into a network
    conformal \( \Delta' \) by setting the additional neurons to zero.
    Thus, \( \calDl_{\Delta} \subseteq \calDl_{\Delta'} \).
    Since
    \( \calDl_{\Delta} \subseteq \cbr{ \emph{\diag}(z) : z \in \cbr{0, 1}^n} \),
    any monotone increasing sequence \( \Delta_{k} \) defines a bounded
    monotone increasing sequence \( \calDl_{\Delta_k} \) which must have
    a limit.
    If every entry of \( \Delta_k \) diverges as \( k \rightarrow \infty \),
    then this limit is the same no matter the choice of sequence.
    We call it \( \calDl \).
\end{proof}

\layerElimination*
\begin{proof}
    We start by defining a suitable \( T^{(l+1)} \) as follows:
    \begin{equation}\label{eq:t-plus-characterization}
        T^{(l+1)}_{j_1, \ldots, j_l}
        = \sum_{i_l \in \calI^{(l)}_{j_l}}
        T^{(l)}_{j_1, \ldots, j_{l-1}, i_l} \otimes W^{(l+1)}_{i_l}
        \text{ where } \calI^{(l)}_{j_{l}}
        = \cbr{ i_l : T^{(l)}_{i_l} \in \calK^{(l)}_{j_{l}}}
        ,
    \end{equation}
    where we recall that the index \( j_l \) ranges over the activation
    patterns achievable by the \( l^{\text{th}} \) layer,
    \[
        \calD^{(l)} = \cbr{
            D^{(l)}_{j_l}
            = \text{diag}\rbr{\mathds{1}(X^{(l)} \odot T^{(l)}_{i_l} \geq 0) }
            : T^{(l)} \in \calG^{(l)}, \, i_l \in [d_l]
        }.
    \]
    Remember that \( p_l = \abs{\calD^{(l)}} \) denotes the cardinality of this set.
    Using this notation, we can re-index the sum in
    \cref{eq:layer-to-tensorize} by activation patterns as follows:
    \begin{align*}
        Z^{(l+2)}
         & = \rbr{\sum_{i_l=1}^{d_l} \rbr{X^{(l)} \odot T^{(l)}_{i_l}}_+
        W^{(l+1)}_{i_l}}_+                                               \\
         & = \rbr{\sum_{j_l=1}^{p_l} \sum_{i_l \in \calI^{(l)}_{j_l}}
            \rbr{X^{(l)} \odot T^{(l)}_{i_l}}_+
            W^{(l+1)}_{i_l}}_+.
        \intertext{Linearizing the activation pattern now gives,}
        Z^{(l+2)}
         & = \rbr{\sum_{j_l=1}^{p_l}
            \sum_{i_l \in \calI^{(l)}_{j_l}}
            D^{(l)}_{j_l} \sbr{X^{(l)} \odot T^{(l)}_{i_l}}
        W^{(l+1)}_{m, \cdot}}_+                                          \\
         & = \rbr{\sum_{j_l=1}^{p_l}
        \sum_{j_1, \ldots, j_{l-1}} D^{(l)}_{j_l} X^{(l)}_{j_1, \ldots j_{l-1}}
        \rbr{
        \sum_{i_l \in \calI^{(l)}_{j_l}}
        T^{(l)}_{j_1, \ldots, j_{l-1}, i_l}
        \otimes W^{(l+1)}_{i_l}}}_+                                      \\
         & = \rbr{X^{(l+1)} \odot T^{(l+1)}}_+.
    \end{align*}
    It should be clear from the way that we constructed \( T^{(l+1)} \) that
    it is contained \( \calG^{(l+1)} \).
    This establishes one direction of theorem.
    The other direction can be shown by instead assuming that
    \( T^{(l+1)} \in \calG^{(l+1)} \) and then noting that the characterization
    in \cref{eq:t-plus-characterization} immediately follows for some
    \( T^{(l)} \in \calG^{(l)} \) and \( W^{(l+1)} \in \R^{d_l \times d_{l+1}} \).
    The preceding calculation is then reversed to obtain
    \cref{eq:layer-to-tensorize}.
\end{proof}

